\documentclass{article}

\usepackage[utf8]{inputenc}     % pdfLaTeX
\usepackage[T1]{fontenc}
\usepackage[magyar]{babel}

% Set page size and margins
% Replace `letterpaper' with `a4paper' for UK/EU standard size
\usepackage[a4paper,top=2cm,bottom=2cm,left=3cm,right=3cm,marginparwidth=1.75cm]{geometry}
%for arrows
\usepackage{textcomp}

\title{Névmások}
\author{Lengyel Zoltán Péter \thanks{Csukás Ibolya tanítása alapján}}
\date{Legutóbb frissítve: 2026. Március 15.}

\begin{document}

\maketitle
\tableofcontents
\newpage

\begin{center}
    {\Huge Névmások}
\end{center}

\begin{itemize}
    \item Helyettesítő szavak
    \item Helyettesíthetnek: \begin{itemize}
        \item Főnevet
        \item Melléknevet
        \item Számnevet
        \item Határozószót
    \end{itemize}
\end{itemize}

\section{Főnevet helyettesítő névmások}
\subsection{Személyes névmás}
\begin{itemize}
    \item én, te, ő, mi, ti, ők, Ön, Maga, Kegyed
    \item Ragos alakok pl: \begin{itemize}
        \item engem, téged, őt, minket, titeket, őket, Önt, Magát, Kegyedet
        \item velem, veled, vele, velünk, veletek, velük, Önnel, Magával, Kegyeddel
        \item rólam, rólad, róla, rólunk, rólatok, róluk, Önről, Magáról, Kegyedről
        \item nálam, nálad, nála, nálunk, nálatok, náluk, Önnél, Magánál, Kegyednél
    \end{itemize}
    \item Névutós alakok pl: \begin{itemize}
        \item miattam, miattad, miatta, miattunk, miattatok, miattuk
        \item mellettem, melletted, mellette, mellettünk, mellettetek, mellettük
        \item alattam, alattad, alatta, alattunk, alattatok, alattuk
    \end{itemize}
\end{itemize}

\subsection{Birtokos névmás}
\begin{itemize}
    \item kifejezik a birtokos számát, személyét és a birtok számát
    \\ \begin{tabular}{|l|l|l|}
\hline
Birtokos & Egy birtok & Több birtok \\ \hline
E/1 & enyém & enyéim \\ \hline
E/2 & tied & tieid  \\ \hline
E/3 & övé & övéi \\ \hline
T/1 & miénk  & mieink \\ \hline
T/2 & tietek & tieitek \\ \hline
T/3 & övék & övéik \\ \hline
\end{tabular}
\end{itemize}

\subsection{Visszaható névmás}
\begin{itemize}
    \item maga + toldalék
    \item Kifejezheti: \begin{itemize}
        \item Egyedüllétet pl: A gyerekek egész nap maguk voltak.
        \item Egyedüli, segítség nélküli cselekvést pl: Péter maga készítette el az ebédet. 
    \end{itemize}
\end{itemize}

\subsection{Kölcsönös névmás}
\begin{itemize}
    \item egymás + toldalék
    \item pl. egymást, egymásnak
\end{itemize}

\newpage
\section{Főnevet, melléknevet, számnevet vagy határozószót helyettesítő névmások}
\begin{tabular}{|l|l|l|l|l|}
\hline
     Névmás & Főnévi & Melléknévi & Számnévi & Határozószói \\ \hline
     Kérdő névmás & ki? & milyen? & mennyi? & hol? \\
     & mi? & mekkora? & hányadik? & hogyan? \\ 
     & & & & mikor? \\ \hline
     Mutató névmás & ez - az & ilyen - olyan & ennyi - annyi & itt - ott \\
     & ugyanez - ugyanaz & amilyen - amolyan & emennyi - amannyi & emitt - amott \\ 
     & & & & így - úgy \\
     & & & & ekkor - akkor \\ \hline
     Határozatlan névmás & valaki & valamilyen & valamennyi & valahol \\
     & egyik - másik & valahanyadik & néhol & \\ \hline
     Általános névmás & mindenki & & & mindenhol \\ 
     & minden & & & \\
     & bárki & bármilyen & bármennyi & bárhol \\
     & bármi & bármekkora  &  & bárhogyan \\
     & & & & bármikor \\ 
     & akárki & akármilyen & akármennyi & akárhogy \\
     & akármi & akármekkora & & akármikor \\
     & senki & & semennyi & sehol \\
     & semmi & & & sehogy \\ \hline
     Vonatkozó névmás & aki & amilyen & amennyi & ahol \\
     & ami & amekkora & ahány & ahogy \\
     & & & & amikor \\ \hline
     
\end{tabular}

\end{document}